\documentclass{article}
\usepackage{amsmath,amssymb,amsthm}
\usepackage{algorithm}
\usepackage{algpseudocode}
\usepackage{hyperref}
\newtheorem{theorem}{Theorem}
\newtheorem{lemma}{Lemma}
\newtheorem{assumption}{Assumption}
\begin{document}

\section*{Trust Region Gradient Method (TRGM)}

\section*{Mathematical Formulation}
Let $f:\mathbb{R}^n\rightarrow\mathbb{R}$ be convex and $L$‑smooth, i.e.
\[
\|\nabla f(x)-\nabla f(y)\|\le L\|x-y\|\qquad\forall x,y\in\mathbb{R}^n .
\]
At iteration $k$ we build the quadratic model
\[
m_k(p)=f(x_k)+\langle\nabla f(x_k),p\rangle+\frac{L}{2}\|p\|^2 ,
\]
and we (approximately) solve the trust–region subproblem
\begin{equation}
\label{eq:subprob}
\min_{p\in\mathbb{R}^n}\; m_k(p)\quad\text{s.t.}\quad \|p\|\le \Delta_k ,
\end{equation}
where $\Delta_k>0$ is the trust‑region radius.  
A closed‑form approximate solution that is sufficient for the analysis is
\[
p_k=
-\min\!\Bigl\{\Delta_k,\;\frac{\|\nabla f(x_k)\|}{L}\Bigr\}\,
\frac{\nabla f(x_k)}{\|\nabla f(x_k)\|}\qquad
\bigl(\text{if }\nabla f(x_k)\neq0\bigr) .
\tag{TRGM‑step}
\]
The new iterate is
\[
x_{k+1}=x_k+p_k .
\]

\section*{Pseudocode}
\begin{algorithm}[H]
\caption{Trust Region Gradient Method (TRGM)}\label{alg:trgm}
\begin{algorithmic}[1]
\Require Initial point $x_0\in\mathbb{R}^n$, initial radius $\Delta_0>0$, Lipschitz constant $L$, 
        radius update parameters $0<\gamma_{\text{inc}}<\gamma_{\text{dec}}<1$.
\For{$k=0,1,2,\dots$}
    \State Compute gradient $g_k=\nabla f(x_k)$.
    \If{$\|g_k\|=0$}
        \State \textbf{stop} (optimal point reached).
    \EndIf
    \State Compute step
    \[
    p_k=-\min\!\Bigl\{\Delta_k,\;\frac{\|g_k\|}{L}\Bigr\}\,\frac{g_k}{\|g_k\|}.
    \]
    \State Evaluate trial point $x_k^{\text{trial}}=x_k+p_k$ and actual reduction
    $ \rho_k=\dfrac{f(x_k)-f(x_k^{\text{trial}})}{m_k(0)-m_k(p_k)}$.
    \If{$\rho_k\ge \eta$}   \Comment{accept step}
        \State $x_{k+1}\gets x_k^{\text{trial}}$.
    \Else
        \State $x_{k+1}\gets x_k$.
    \EndIf
    \State Update radius
    \[
    \Delta_{k+1}\gets 
    \begin{cases}
    \min\{\Delta_{\max},\;\gamma_{\text{inc}}\Delta_k\}, & \rho_k\ge \eta,\\[4pt]
    \gamma_{\text{dec}}\Delta_k, & \rho_k<\eta .
    \end{cases}
    \]
    \If{termination criterion satisfied (e.g., $\|g_k\|\le\varepsilon$)}\State \textbf{break}\EndIf
\EndFor
\end{algorithmic}
\end{algorithm}

\section*{Dry Run Example}
Consider the one–dimensional convex smooth function
\[
f(w)=(w-3)^2 .
\]
We have $\nabla f(w)=2(w-3)$ and $L=2$ (since $f''(w)=2$).  
Choose $x_0=0$, initial radius $\Delta_0=1$, and set $\eta=0.1$, $\gamma_{\text{inc}}=1.5$, $\gamma_{\text{dec}}=0.5$, $\Delta_{\max}=10$.

\medskip
\noindent\textbf{Iteration 0}
\[
g_0=\nabla f(0)=2(0-3)=-6,\qquad \|g_0\|=6.
\]
Step size:
\[
\alpha_0=\min\Bigl\{\Delta_0,\frac{\|g_0\|}{L}\Bigr\}
       =\min\{1,\;6/2\}=1.
\]
\[
p_0=-\alpha_0\frac{g_0}{\|g_0\|}= -1\cdot\frac{-6}{6}=1.
\]
Trial point $x_0^{\text{trial}}=0+1=1$.
\[
f(0)=9,\qquad f(1)=4.
\]
Model reduction:
\[
m_0(0)-m_0(p_0)=\langle g_0,p_0\rangle+\frac{L}{2}\|p_0\|^2
               =(-6)(1)+\frac{2}{2}(1)= -6+1=-5.
\]
Actual reduction $=9-4=5$, thus
\[
\rho_0=\frac{5}{5}=1\ge\eta .
\]
Accept step: $x_1=1$, increase radius: $\Delta_1=\min\{10,1.5\}=1.5$.

\medskip
\noindent\textbf{Iteration 1}
\[
g_1=\nabla f(1)=2(1-3)=-4,\;\|g_1\|=4.
\]
\[
\alpha_1=\min\{1.5,\;4/2\}=1.5.
\]
\[
p_1=-1.5\frac{-4}{4}=1.5,\qquad x_1^{\text{trial}}=1+1.5=2.5.
\]
\[
f(1)=4,\; f(2.5)=(2.5-3)^2=0.25.
\]
Model reduction:
\[
\langle g_1,p_1\rangle+\frac{L}{2}\|p_1\|^2=(-4)(1.5)+\frac{2}{2}(1.5)^2=-6+2.25=-3.75.
\]
Actual reduction $=4-0.25=3.75$, hence $\rho_1=1\ge\eta$.
Accept step: $x_2=2.5$, increase radius: $\Delta_2=\min\{10,1.5\cdot1.5\}=2.25$.

\medskip
\noindent\textbf{Iteration 2}
\[
g_2=\nabla f(2.5)=2(2.5-3)=-1,\;\|g_2\|=1.
\]
\[
\alpha_2=\min\{2.25,\;1/2\}=0.5.
\]
\[
p_2=-0.5\frac{-1}{1}=0.5,\qquad x_2^{\text{trial}}=2.5+0.5=3.0.
\]
\[
f(2.5)=0.25,\; f(3)=0.
\]
Model reduction:
\[
\langle g_2,p_2\rangle+\frac{L}{2}\|p_2\|^2=(-1)(0.5)+\frac{2}{2}(0.5)^2=-0.5+0.25=-0.25.
\]
Actual reduction $=0.25$, $\rho_2=1\ge\eta$.
Accept step: $x_3=3$, radius would increase but algorithm stops because $\|g_2\|=0$ (optimal point).

The sequence $\{x_k\}$ rapidly converges to the minimiser $w^\star=3$.

\newpage
\section*{Assumptions}
\begin{assumption}[Convexity and Smoothness]\label{ass:convexsmooth}
The objective $f:\mathbb{R}^n\to\mathbb{R}$ is convex and $L$‑smooth:
\begin{enumerate}
\item \textbf{Convexity:} $f(y)\ge f(x)+\langle\nabla f(x),y-x\rangle,\;\forall x,y$.
\item \textbf{Lipschitz gradient:} $\|\nabla f(x)-\nabla f(y)\|\le L\|x-y\|,\;\forall x,y$.
\end{enumerate}
\end{assumption}
\begin{assumption}[Trust‑Region Radius]\label{ass:radius}
The radii satisfy $\Delta_k\ge \frac{\|\nabla f(x_k)\|}{L}$ for all $k$ (this can be enforced by a simple back‑tracking rule). Moreover $\Delta_k\le\Delta_{\max}<\infty$.
\end{assumption}
\begin{assumption}[Exact Model]\label{ass:model}
The model $m_k$ defined in \eqref{eq:subprob} uses the exact Lipschitz constant $L$; hence $m_k(p)\ge f(x_k+p)$ for all $\|p\|\le\Delta_k$.
\end{assumption}

\section*{Main Convergence Theorem}
\begin{theorem}[Sublinear $O(1/k)$ rate]\label{thm:rate}
Let Assumptions \ref{ass:convexsmooth}–\ref{ass:model} hold and let $\{x_k\}$ be the iterates generated by TRGM with the step \eqref{TRGM‑step}. Denote $x^\star\in\arg\min f$ and $R:=\|x_0-x^\star\|$. Then for all $K\ge1$,
\[
f(x_K)-f(x^\star)\;\le\;\frac{2L R^{2}}{K}\, .
\]
Consequently $\displaystyle\lim_{k\to\infty}f(x_k)=f(x^\star)$ and $\|x_k-x^\star\|\to0$.
\end{theorem}

\section*{Theoretical Properties}
\begin{itemize}
\item \textbf{Monotonicity.} Under Assumption \ref{ass:model} the actual reduction is at least the model reduction, guaranteeing $f(x_{k+1})\le f(x_k)$.
\item \textbf{Stability w.r.t.\ $\Delta_k$.} If $\Delta_k$ is chosen larger than $\|\nabla f(x_k)\|/L$, the step $p_k$ coincides with the exact Cauchy point of the model, which yields maximal possible decrease for the quadratic model.
\item \textbf{Comparison to Gradient Descent.} When $\Delta_k$ is taken unbounded the method reduces to the standard gradient step with step size $1/L$, reproducing the classical $O(1/k)$ rate. The trust region therefore provides an adaptive safeguard without sacrificing the rate.
\item \textbf{Hyperparameter Sensitivity.} The constants $\gamma_{\text{inc}},\gamma_{\text{dec}}$ affect only the radius evolution; the $O(1/k)$ bound holds for any choice as long as Assumption \ref{ass:radius} is respected.
\end{itemize}

\section*{Discussion}
The theorem shows that, even with only an approximate solution of the subproblem (the Cauchy point), the trust‑region scheme attains the optimal sublinear rate for convex smooth optimisation. The proof hinges on the exact quadratic model and the radius condition guaranteeing that the step length is at least $\|\nabla f(x_k)\|/L$, which reproduces the descent obtained by the classical $1/L$ gradient step. Hence TRGM can be viewed as a robust variant of gradient descent that automatically adjusts step lengths via the trust‑region radius.

\newpage
\section*{Preliminaries (Definitions and Lemmas)}
\begin{lemma}[Smoothness inequality]\label{lem:smooth}
If $f$ is $L$‑smooth then for all $x,p$,
\[
f(x+p)\le f(x)+\langle\nabla f(x),p\rangle+\frac{L}{2}\|p\|^{2}.
\]
\end{lemma}
\begin{proof}
Standard consequence of the integral form of the gradient Lipschitz condition. \qedhere
\end{proof}

\begin{lemma}[Convexity inequality]\label{lem:convex}
For convex $f$,
\[
f(x)-f(x^\star)\le\langle\nabla f(x),x-x^\star\rangle .
\]
\end{lemma}
\begin{proof}
Set $y=x^\star$ in the definition of convexity. \qedhere
\end{proof}

\section*{Main Proof (Step‑by‑Step Derivation)}
We prove Theorem \ref{thm:rate}. Throughout the proof $x^\star$ denotes any optimal solution.

\noindent\textbf{Step 1 (Model decrease).}  
From the definition of $p_k$ and Assumption \ref{ass:radius},
\[
\|p_k\|=\frac{\|\nabla f(x_k)\|}{L}. \tag{1}
\]
Using Lemma \ref{lem:smooth} with $p=p_k$ and the expression of $m_k$,
\[
\begin{aligned}
f(x_k)-f(x_{k+1})
&\ge f(x_k)-\bigl[f(x_k)+\langle g_k,p_k\rangle+\tfrac{L}{2}\|p_k\|^{2}\bigr]   \\
&= -\langle g_k,p_k\rangle-\tfrac{L}{2}\|p_k\|^{2}. \tag{2}
\end{aligned}
\]
Insert $p_k=-\frac{\|g_k\|}{L}\frac{g_k}{\|g_k\|}$ from (1):
\[
-\langle g_k,p_k\rangle = \frac{\|g_k\|^{2}}{L},\qquad
\|p_k\|^{2}= \frac{\|g_k\|^{2}}{L^{2}} .
\]
Hence
\[
f(x_k)-f(x_{k+1})\ge \frac{\|g_k\|^{2}}{L}-\frac{L}{2}\frac{\|g_k\|^{2}}{L^{2}}
= \frac{\|g_k\|^{2}}{2L}. \tag{3}
\]
\noindent\textbf{[Step 1]}

\noindent\textbf{Step 2 (Relating gradient norm to suboptimality).}  
From Lemma \ref{lem:convex} with $x=x_k$,
\[
f(x_k)-f(x^\star)\le \langle g_k,x_k-x^\star\rangle
\le \|g_k\|\,\|x_k-x^\star\|. \tag{4}
\]
\noindent\textbf{[Step 2]}

\noindent\textbf{Step 3 (Bounding distance to optimum).}  
Square both sides of (4) and use $ab\le\frac{a^{2}+b^{2}}{2}$:
\[
\bigl(f(x_k)-f(x^\star)\bigr)^{2}
\le \|g_k\|^{2}\|x_k-x^\star\|^{2}
\le \frac{\|g_k\|^{2}}{2}\bigl(\|g_k\|^{2}+ \|x_k-x^\star\|^{2}\bigr). \tag{5}
\]
Because $f$ is convex, the sequence $\{f(x_k)\}$ is non‑increasing, thus $f(x_k)-f(x^\star)\le f(x_0)-f(x^\star)$. Hence $\|g_k\|$ is bounded; we can drop the $\|g_k\|^{2}$ term in (5) and obtain
\[
\bigl(f(x_k)-f(x^\star)\bigr)^{2}
\le \|g_k\|^{2}\,R^{2},\qquad R:=\|x_0-x^\star\|. \tag{6}
\]
\noindent\textbf{[Step 3]}

\noindent\textbf{Step 4 (Recursive inequality).}  
Combine (3) and (6):
\[
f(x_k)-f(x_{k+1})
\ge \frac{1}{2L}\|g_k\|^{2}
\ge \frac{1}{2L}\frac{\bigl(f(x_k)-f(x^\star)\bigr)^{2}}{R^{2}}. \tag{7}
\]
Define $\Delta_k:=f(x_k)-f(x^\star)\ge0$. Then (7) reads
\[
\Delta_k-\Delta_{k+1}\ge \frac{\Delta_k^{2}}{2LR^{2}}. \tag{8}
\]
\noindent\textbf{[Step 4]}

\noindent\textbf{Step 5 (Summation).}  
Rewrite (8) as
\[
\frac{1}{\Delta_{k+1}}-\frac{1}{\Delta_{k}}
\le -\frac{1}{2LR^{2}}. \tag{9}
\]
Summing (9) for $k=0,\dots,K-1$ yields
\[
\frac{1}{\Delta_{K}}-\frac{1}{\Delta_{0}}
\le -\frac{K}{2LR^{2}}.
\]
Since $\Delta_{0}=f(x_0)-f(x^\star)\le \frac{L}{2}R^{2}$ (by smoothness of $f$ around $x^\star$), we obtain
\[
\frac{1}{\Delta_{K}}\ge \frac{K}{2LR^{2}}.
\]
Inverting gives
\[
\Delta_{K}\le \frac{2LR^{2}}{K}. \tag{10}
\]
\noindent\textbf{[Step 5]}

\noindent\textbf{Conclusion.}  
Inequality (10) is exactly the claim of Theorem \ref{thm:rate}. $\square$

\section*{Rate Derivation (With Constants)}
The bound \eqref{10} can be written more explicitly as
\[
f(x_K)-f(x^\star)\le \underbrace{\frac{2L}{K}}_{\text{iteration factor}}\,
\underbrace{\|x_0-x^\star\|^{2}}_{\text{initial distance}} .
\]
Thus the convergence constant is $2L$, identical to that of the classic gradient method with step size $1/L$.

\section*{Special Cases}
\begin{itemize}
\item \textbf{Exact solution of \eqref{eq:subprob}.}  
If the subproblem is solved exactly, the step satisfies
\[
p_k^{\star}= -\frac{1}{L}\nabla f(x_k)\quad\text{whenever }\|p_k^{\star}\|\le\Delta_k,
\]
and the same proof applies, giving the same $O(1/k)$ bound.
\item \textbf{Strongly convex $f$.}  
If $f$ is $\mu$‑strongly convex, Lemma \ref{lem:convex} can be strengthened to
\[
f(x_k)-f(x^\star)\ge \frac{\mu}{2}\|x_k-x^\star\|^{2},
\]
which combined with (3) yields a linear rate
\[
f(x_k)-f(x^\star)\le \bigl(1-\tfrac{\mu}{L}\bigr)^{k}\bigl(f(x_0)-f(x^\star)\bigr).
\]
\end{itemize}

\end{document}